\section{Introduction}
\label{sec:intro}

%-------------------------------------------------------------------------
Deep convolutional neural networks have led to a series of breakthroughs for image classification. Deep networks naturally integrate low/mid/high-level features and classifiers in an end-to-end multilayer fashion, and the "levels" of features can be enriched by the number of stacked layers(layers). Recent evidence reveals that network depth is of crucial importance, and leading results on the challenging ImageNet dataset all exploit\words{(采用v.)} "very deep" models, with a depth of sixteen to thirty. Many other nontrivial visual recognition tasks have also greatly benefited from very deep models.\ygm{(网路的深度是影响视觉识别相关任务的关键,加深网络层数可以学习到图像中高层次的视觉特征,从未达到更好的分类效果)}

Driven by the significance of \textbf{depth}, a question arises: \textit{Is learning better networks as easy as stacking more layers?}\ygm{(在认识到了层数的重要性后提出问题：学习更好的网络是否就像堆叠更多的层一样简单?)} An obstacle to answering this question was the notorious problem of vanishing/exploding gradients, which hamper convergence from the beginning. This problem, however, has been largely \textbf{addressed}\words{(已解决v.)} by normalized initialization and intermediate normalization layers, which enable networks with tens of layers to start converging for stochastic gradient descent(SGD) with backpropagation.

When deeper networks are able to start converging, a \textit{degradation}\words{(性能退化adj.)} problem has been exposed: with the network depth increasing, accuracy gets saturated (which might be unsurprising) and then degrades rapidly.\ygm{(随机梯度的反向传播训练可以使深层的网络收敛,但是随着网络层数加深,准确率会达到饱和并迅速退化)} Unexpectedly, such \textbf{degradation is not caused by overfitting}, and \textbf{adding more layers to a suitably deep models leads to higher training error}, as reported in [10,41] and thoroughly verified by our experiments, Fig.1 shows a typical example. \ygm{(对于一个深度适中的网络,增加层数会带来更高的训练误差[文献10、41以及文本的实验中证实了这个问题])}

The degradation(\textbf{of training accuracy}) indicates that not all systems are similarly easy to optimize.\ygm{(如果是过拟合,那么训练误差应该仍然很好;训练误差本身变差,说明问题出在优化上,网络深了但是SGD优化算法找不到更好的解)} Let us consider a shallower architecture and its deeper counterpart that adds more layers onto it. There exists a solution by construction to the deeper model: the added layers are identity mapping, and the other layers are copied from the learned shallower model.\ygm{(对于深层网络而言,存在一个构造解:把浅网络学好的权重原样拷进前几层，新增的层设为恒等映射)} The existence of this constructed solution indicates that a deeper model should produce no higher training error than its shallower counterpart.\ygm{(既然这样的解存在，深网络的训练误差理应不高于浅网络)} But experiments show that our current solvers on hand unable to find solutions that are comparably good or better than the constructed solution (or unable to do so in feasible time).\ygm{(但实验表明，我们手头的求解器找不到与构造解相当或更好的)}

In this paper, we address the degradation problem by introducing a \textit{deep residual learning} framework.\ygm{(提出方法)} Instead of hoping each few stacked layers directly fit a desired underlying mapping, we explicitly let these layers fit a residual mapping.\ygm{(目标转变:拟合底层映射到拟合残差映射)} Formally, denoting the desired underlying mapping as $\mathcal{H}(\mathbf{x})$, we let the stacked nonlinear layers fit another mapping of $\mathcal{F}(\mathbf{x}) := \mathcal{H}(\mathbf{x}) - \mathbf{x}$. The original mapping is recast into $\mathcal{F}(\mathbf{x}) + \mathbf{x}$. \textbf{We hypothesize} that it is easier to optimize the residual mapping than to optimize the original, unreferenced mapping. \textbf{To the extreme},if an identity mapping were optimal, it would be easier to push the residual to zero than to fit an identity mapping by a stack of nonlinear layers.\ygm{(分析了特殊情况下的适用性)}

The formulation of $\mathcal{F}(\mathbf{x}) + \mathbf{x}$ can be realized by feedforward neural networks with "shortcut connections"(Fig.2). Shortcut connections [2,33,48] are those skipping one or more layers. In our case, the shortcut connections simply perform \textit{identity} mapping, and their outputs are added to the outputs of the stacked layers (Fig.2).\ygm{(解释方法的具体实现)} \textbf{Identity shortcut connections add neither extra parameter nor computational complexity.} The entire network can still be trained end-to-end by SGD with backpropagation, and can be easily implemented using common libraries (e.g., Caffe[19]) without modifying the solvers.\ygm{(解释方法的优势)}

We present comprehensive experiments on ImageNet to show the degradation problem and evaluate our method. We show that: 1)Our extremely deep residual nets are \underline{easy to optimize}, but the counterpart "plain" nets (that simply stack layers) exhibit higher training error when the depth increases; 2)Our depth residual nets can easily enjoy \underline{accuracy gains} from greatly increased depth, producing results substantially\words{(显著地adv.)} better than previous networks.\ygm{(容易优化、层数加深不易退化)}

Similar phenomena are also shown on the CIFAR-10 set, suggesting that optimization diffculties and the effects of our method \textbf{are not just akin} a particular dataset.\ygm{(多数据集泛化性证明模型的效果)} We present successfully trained models on this dataset with over 100 layers, and explore models with over 1000 layers.

On the ImageNet classification dataset [35], we obtain excellent results by extremely deep residual nets. Our 152 layers residual net is the deepest network ever presented on ImageNet, while still having lower complexity than VGG nets [40]. Our ensemble has \textbf{3.57\%} top-5 error on the ImageNet \textit{test set}, and \emph{won the 1st place in the ILSVRC 2015 classification competition}.  The extremely deep representations alse have excellent generalization performance on other recognition tasks, and lead us to further \emph{win the 1st places on: ImageNet detection, ImageNet localization, COCO detection, and COCO segmentation} in ILSVRC \& COCO 2015 competitions. This strong evidence shows that the residual learning principle is generic, and we expect that it is applicable in other vision and non-vision problems.\ygm{(通过竞赛效果证明模型的普适性)}


